\section{代码能力与 AI 辅助开发}

\subsection{深度学习框架入门（PyTorch）}
必须掌握至少一种主流框架；本组默认推荐 PyTorch。
\begin{itemize}
  \item 官方 60 分钟速成：\url{https://pytorch.org/tutorials/beginner/deep_learning_60min_blitz.html}
  \item 中文解读参考：\url{https://zhuanlan.zhihu.com/p/86875507}
  \item 大模型相关 PyTorch 入门（视频号 CS336 相关分享）：可在微信视频号检索「斯坦福 CS336 PyTorch」
\end{itemize}
配套练习：李沐《动手学深度学习》B 站课程：\\
\url{https://www.bilibili.com/video/BV1if4y147hS/}

建议路径：官方 tutorial $\rightarrow$ 复现一个小模型（如 MNIST/CIFAR）$\rightarrow$ 阅读领域经典开源实现。

\subsection{推荐的工程目录布局}
阅读或自建训练仓库时，可参考如下结构：
\begin{itemize}
  \item \texttt{config/}：训练超参、路径配置；
  \item \texttt{dataset/} 或 \texttt{dataloader/}：数据读取与增强；
  \item \texttt{models/}：网络模块；
  \item \texttt{metrics/}：损失与评价指标；
  \item \texttt{utils/}：通用工具；
  \item \texttt{scripts/}：训练/评测启动脚本（服务器可用 \texttt{nohup}，集群可用 \texttt{sbatch}）；
  \item \texttt{weights/}： checkpoint；
  \item \texttt{train.py}：主入口。
\end{itemize}

\subsection{如何找论文对应代码}
\begin{enumerate}
  \item 论文摘要/介绍末尾的 GitHub 链接；
  \item 用项目名在 GitHub 搜索；
  \item Papers with Code 按论文名检索；
  \item 必要时礼貌邮件作者（说明身份与用途）。
\end{enumerate}
工业界代码更新快、社区强，但医学场景未必开箱即用；医学专用代码更贴数据，但维护与数据公开性可能较弱。建议各选一份高质量实现精读模仿。

\subsection{读代码的检查清单}
\begin{itemize}
  \item 任务定义：源码要解决什么问题；
  \item 数据：类型、路径、预处理；
  \item 环境：依赖、版本、硬件；
  \item 能否跑通、输出是什么；
  \item 整体架构与关键模块对应论文哪一段。
\end{itemize}

\subsection{Claude Code / AI 辅助编程}
AI 适合辅助写样板代码、解释报错、阅读陌生仓库；但\textbf{不能}替代你对方法与实验的理解，生成内容必须核对。
入门资料：
\begin{itemize}
  \item 完全指南：\url{https://mp.weixin.qq.com/s/5uhudUPNKLYwxZXB0_IH_g}
  \item 命令与工作流：\url{https://mp.weixin.qq.com/s/pwozo7PSu8RGedfzs0E6BA?scene=1}
  \item 背景阅读：\url{https://mp.weixin.qq.com/s/YHDSJayeDwmtqY1pGT9MnA}
\end{itemize}
练习建议：先学会清晰下指令（目标、约束、文件范围），再让 AI 解释/改错，最后自己跑通并写进实验记录。

\subsection{训练与调试习惯}
\begin{itemize}
  \item 每次尽量只改一个变量，便于归因；
  \item 监控损失与关键指标，区分过拟合/欠拟合；
  \item 保存随机种子、配置文件与代码 commit，保证可复现；
  \item 困难时按「数据质量 $\rightarrow$ 预处理 $\rightarrow$ 增强 $\rightarrow$ 结构 $\rightarrow$ 超参 $\rightarrow$ 预训练/迁移」排查。
\end{itemize}
